% ================================================================
% Capitolul 3 -- Stadiu actual
% Lucrare de licenta: Kernel Trap -- Sistem Inteligent de Prevenire
% a Intruziunilor cu Componente de Honeypot si Invatare Automata
% Autor: Banciu George-Horatiu
% ================================================================

Capitolul de față trece în revistă stadiul actual al cercetării în domeniile
care converg în arhitectura sistemului propus, evaluează soluțiile existente
și identifică golul pe care lucrarea de față urmărește să îl acopere. KernelTrap
se situează la intersecția a patru direcții de cercetare active în literatura de
specialitate din ultimii ani: monitorizarea comportamentală la nivel de kernel
prin tehnologia eBPF, detectarea anomaliilor pe baza apelurilor de sistem prin
metode de învățare automată nesupervizată, honeypot-urile cu interacțiune
ridicată pentru analiza atacurilor reale și mecanismele de apărare activă
bazate pe redirecționare transparentă a sesiunilor compromise. Pentru fiecare
direcție sunt prezentate atât lucrările fundamentale, cât și contribuțiile
recente (2022--2025), ale căror rezultate fundamentează deciziile de
proiectare expuse în capitolul următor.

\section{Sisteme HIDS bazate pe eBPF}

Tranziția dinspre modulele de nucleu tradiționale (LKM) către programe eBPF
reprezintă cea mai semnificativă schimbare în domeniul
monitorizării de securitate la nivel de gazdă din ultimul deceniu.
Soldani et al.~\cite{soldani2023ebpf} oferă cea mai cuprinzătoare analiză
recentă a tehnologiei eBPF pentru observabilitate cloud-native și securitate
runtime, demonstrând că suprafața redusă de atac, garanțiile statice oferite de
verificator și costul redus de instrumentare fac din eBPF alegerea naturală
pentru sisteme HIDS moderne.
Hao et al.~\cite{hao2024ebpfruntime} descriu în detaliu implementarea
runtime-ului eBPF în nucleul Linux 6.7, confirmând maturitatea ecosistemului ca
platformă stabilă pentru produse de securitate.

În spațiul soluțiilor open-source, patru proiecte domină piața:
\textbf{Falco} (CNCF, originar de la Sysdig), care folosește eBPF pentru
colectare și un motor de reguli declarative pentru detecție;
\textbf{Tetragon} (Cilium/Isovalent), care diferă prin capacitatea de aplicare
sincronă a politicilor direct în kernel space;
\textbf{KubeArmor} (Accuknox), care exploatează BPF-LSM pentru aplicarea de
politici de securitate;
și \textbf{Tracee} (Aqua Security), care a fost adoptat și de KernelTrap
pentru capacitățile sale de tracing fără configurare prealabilă pe nuclee
diferite (CO:RE/BTF).
Syairozi și Arizal~\cite{syairozi2025ebpfcomparative} efectuează cea mai
recentă comparație empirică a acestor sisteme pe un set de atacuri OWASP
Kubernetes Top~10, măsurând timpul mediu până la detecție (MTTD), rata de
detecție, rata de fals pozitive și consumul de resurse.
Concluzia studiului este că Tracee oferă cea mai bună acoperire a evenimentelor
LSM și cea mai mică amprentă pe CPU dintre soluțiile testate, susținând
alegerea făcută în cadrul prezentei lucrări.

Două lucrări recente sunt deosebit de relevante pentru arhitectura KernelTrap.
eBPF-PATROL~\cite{ebpfpatrol2025} propune un agent runtime ușor pentru
interceptarea apelurilor de sistem și aplicarea de politici, raportând un
overhead sub 2,5\% un punct de referință apropiat de obiectivele de
performanță ale agentului din KernelTrap.
CryptoGuard~\cite{cryptoguard2024}, publicat la DSN 2024, combină monitorizarea
cu fereastră glisantă a syscall-urilor cu o rețea neuronală evaluată direct în
kernel, atingând un overhead remarcabil de 0,06\% pe CPU.
Aceste lucrări demonstrează fezabilitatea sistemelor HIDS de tip
\textit{always-on} în medii de producție, însă niciuna nu integrează un
mecanism de răspuns activ comparabil cu pivotarea propusă de KernelTrap.

O contribuție notabilă este lucrarea fundamentală a lui Fournier et
al.~\cite{fournier2021ebpf} de la Datadog, care discută în detaliu
compromisurile dintre kprobe-uri, tracepoint-uri și hook-urile LSM, precum și
problemele de tip TOCTOU (\textit{Time-of-Check to Time-of-Use})~ aspecte
direct relevante pentru proiectarea agentului KernelTrap.

\section{Învățare automată nesupervizată pentru detectarea intruziunilor}

Limitările abordărilor supervizate aplicate la detectarea intruziunilor sunt
bine documentate în literatură.
Zoppi și Ceccarelli~\cite{zoppi2022unsupervised,zoppi2022stacking} argumentează 
că modelele bazate pe semnături și clasificatorii binari
supervizați eșuează sistematic la detectarea atacurilor zero-day, întrucât
distribuția datelor de atac din timpul antrenării nu reflectă varietatea
atacurilor întâlnite în producție.
Studiul lor de pe TON\_IoT și CICIDS demonstrează că un detector nesupervizat
(în particular Isolation Forest) păstrează o sensibilitate reală pe atacuri
necunoscute, în timp ce clasificatorii supervizați colapsează la performanțe
sub nivelul aleatoriu pe această clasă de scenarii.

Această observație a generat o linie consistentă de cercetare bazată pe
Isolation Forest~\cite{liu2008iforest} și variante ale acestuia.
Studiul comparativ publicat în 2025 în jurnalul
\textit{Computers}~\cite{oneclass2025survey} testează direct OCSVM, Autoencoder,
VAE și Isolation Forest pe seturi de date industriale, raportând că IF oferă
cel mai bun compromis între acuratețe, complexitate computațională și
interpretabilitate.
Anand et al.~\cite{anand2025iot} confirmă rezultatul pe TON\_IoT, demonstrând
superioritatea IF față de OCSVM pe date de înaltă dimensionalitate cu
dezechilibre puternice între clase.

În domeniul specific al apelurilor de sistem, cercetarea modernă se grupează în
jurul setului de date BETH~\cite{highnam2021beth}, introdus de Highnam et al.
în 2021 ca prima sursă publică de evenimente de syscall colectate dintr-un
scenariu real de honeypot în mediu cloud.
Lakha et al.~\cite{lakha2022beth} publică în 2022 una dintre cele mai citate
lucrări bazate pe BETH, folosind o combinație de Graph Neural Network și
Transformer pentru detectarea anomaliilor; rezultatele lor
($\text{AUROC} > 0{,}90$) reprezintă un punct de referință competitiv, însă cu
costuri computaționale semnificativ mai mari decât IF.
Highnam et al.~\cite{highnam2023adaptive} revin în 2023 cu o lucrare despre
proiectarea experimentală adaptivă pentru colectarea de date de intruziune,
semnalând că BETH continuă să fie întreținut activ ca resursă academică.

Cea mai apropiată lucrare metodologică de KernelTrap o reprezintă studiul lui
Atawneh et al.~\cite{atawneh2025symmetry}, publicat în 2025 în
\textit{Symmetry}, care aplică Isolation Forest direct pe setul BETH și
depășește performanțele de referință raportate în lucrarea originală a setului
de date.
Diferența esențială dintre \cite{atawneh2025symmetry} și KernelTrap este că
primul rămâne un studiu de evaluare pe date statice, în timp ce KernelTrap
propune un sistem complet de producție.
Eremin~\cite{eremin2025streams} extinde această linie cu o evaluare a zece
algoritmi de învățare pe fluxuri (incluzând variante de IF) pe setul BETH,
confirmând fezabilitatea detecției nesupervizate în regim de streaming.
Khan et al.~\cite{khan2025marl} explorează o direcție alternativă bazată pe
Multi-Agent Reinforcement Learning aplicat tot pe BETH; deși rezultatele lor
sunt promițătoare, complexitatea de antrenare și interpretabilitatea redusă fac
abordarea improprie pentru un sistem de producție cu cerințe stricte de latență.

În ceea ce privește selecția caracteristicilor pentru ML aplicat la
syscall-uri, Mvula et al.~\cite{mvula2023embeddings} compară diferite tehnici
de extracție a caracteristicilor (one-hot, n-grame, word embeddings) pe
trace-uri de syscall, iar Užupis et al.~\cite{uzupis2024riskbased} introduc o
metodă de grupare bazată pe risc care îmbunătățește acuratețea SVM cu 23,4\%.
Aceste rezultate susțin direct alegerea făcută în KernelTrap de a extinde setul
clasic de 7 caracteristici BETH cu derivate calculate din câmpurile
\texttt{args}, \texttt{processName}, \texttt{timestamp} și \texttt{threadId}
ignorate de versiunile anterioare ale modelelor BETH, ajungând la un total
de 22 de dimensiuni proiectate să echilibreze puterea discriminativă cu costul
computațional al evaluării.

LIGHT-HIDS~\cite{lighthids2025} reprezintă o lucrare recentă (2025) deosebit
de relevantă, propunând un cadru hibrid Deep SVDD + detecție de noutate pe
secvențe de syscall, cu un câștig de viteză la inferență de până la 75 de ori.
Lucrarea confirmă fezabilitatea sistemelor ML-HIDS ușoare în producție, dar nu
integrează nici răspuns activ, nici componente de honeypot.

\section{Praguri adaptive și agregare temporală}

Două direcții complementare sprijină arhitectura de decizie a KernelTrap.
Cui et al.~\cite{cui2022ueba} și Cao et al.~\cite{cao2023ueba} documentează
fundamentarea teoretică a abordării
\textbf{User and Entity Behavior Analytics (UEBA)}, demonstrând că pragurile
per-entitate (calibrate individual pentru fiecare utilizator) reduc semnificativ
rata de fals pozitive comparativ cu un prag global unic~--- o observație
implementată direct în KernelTrap prin pragurile de severitate calibrate
per-UID pe partiția de validare a setului BETH.
Dani et al.~\cite{dani2016adaptive} formalizează utilizarea pragurilor adaptive
bazate pe percentile, justificând alegerea percentilelor 2,0\% și 0,2\% pentru
cele două niveluri de severitate.

Pe direcția agregării temporale, sondajul publicat în
\textit{PeerJ Computer Science}~\cite{streamingsurvey2025} sintetizează
tehnicile de fereastră glisantă, aplicate la detectarea anomaliilor pe fluxuri de
date de securitate.
Algoritmul iForestASD~\cite{ding2024iforestasd} reprezintă una dintre primele
formulări ale combinației Isolation Forest + fereastră glisantă în context de
streaming, iar Yunus et al.~\cite{yunus2024sliding} extind ideea cu o serie de
ferestre paralele care reduc rata de fals pozitive prin votare.
Aceste contribuții susțin direct mecanismul de agregare temporală implementat
în KernelTrap, care combină un prag absolut (cel puțin 5 evenimente de
severitate 2 într-o fereastră de 60 de secunde) cu un prag de rată (cel puțin
30\% din totalul evenimentelor utilizatorului în fereastră să fie de severitate
2), reducând astfel rata de fals pozitive pe utilizatori legitimi cu activitate
intensă.

\section{Honeypot-uri SSH cu interacțiune ridicată}

Honeypot-ul Cowrie rămâne, în literatura recentă, soluția de referință pentru
capturarea atacurilor SSH.
Lucrarea Q-Cowrie~\cite{qcowrie2026}, publicată în 2026 în
\textit{International Journal of Information Security}, propune o variantă
adaptivă bazată pe Q-learning, demonstrând relevanța continuă a Cowrie ca
platformă experimentală.
O direcție inovatoare este reprezentată de Sladić et al.~\cite{sladic2024llm}
în lucrarea \textit{LLM in the Shell} (IEEE EuroS\&PW 2024), care integrează
modele lingvistice mari peste arhitectura Cowrie pentru a genera răspunsuri
shell credibile la comenzi neașteptate ale atacatorului.
Reworr și Volkov~\cite{reworr2024llmhoneypot} aplică o abordare similară
pentru a monitoriza specific agenți de hacking bazați pe LLM, raportând
capturarea unor modele comportamentale distincte față de atacatorii umani.

Pe direcția izolării prin containere, Iwendi et al.~\cite{iwendi2023container}
documentează un sistem honeypot complet containerizat în cloud, iar
HoneyFactory~\cite{honeyfactory2024} propune o arhitectură honeynet bazată pe
containere Docker cu suport pentru orchestrare automată.
Studiul comparativ publicat în 2025 în
\textit{Informatics}~\cite{honeypotsurvey2025} analizează șapte honeypot-uri
contemporane și concluzionează că, pentru scenariul SSH cu interacțiune ridicată,
Cowrie oferă cel mai bun raport între bogăția datelor capturate și costul
operațional.
Aceste lucrări susțin direct alegerea făcută în KernelTrap, unde containerul
honeypot expune un mediu Linux real izolat prin Docker, evitând complexitatea
unei mașini virtuale dedicate.

O limitare comună a tuturor acestor sisteme este însă natura lor pasivă:
atacatorul trebuie să decidă să interacționeze cu honeypot-ul.
Atunci când un atacator a compromis deja o gazdă legitimă printr-un cont SSH
valid, niciunul dintre honeypot-urile descrise mai sus nu îl poate redirecționa
acolo.

\section{Apărare activă și pivotare transparentă}

Lucrarea cea mai apropiată conceptual de mecanismul de pivotare al KernelTrap
este \textit{Cyber Deception Reactive: TCP Stealth Redirection to On-Demand
Honeypots} a lui Lopez et al.~\cite{lopez2024tcpredirection}, publicată în
februarie 2024.
Autorii propun un sistem care interceptează la nivel TCP conexiunile
atacatorului și le redirecționează discret către un honeypot clonat, fără ca
atacatorul să sesizeze tranziția.
Diferențele de concepție față de KernelTrap sunt totuși semnificative: sistemul
lui Lopez et al. operează pe baza unor reguli statice de detecție la nivel de
rețea, în timp ce KernelTrap declanșează pivotarea pe baza unei decizii ML
calibrate pe comportamentul individual al utilizatorului, observat la nivel de
syscall, nu de pachet de rețea.
Mai mult, KernelTrap utilizează semnalul POSIX \texttt{SIGUSR1} combinat cu
reguli \texttt{iptables} pentru a forța reconectarea sesiunii SSH existente, în
timp ce abordarea bazată exclusiv pe NAT a lui Lopez et al. nu poate trata
sesiunile deja stabilite.

Sondajul publicat în \textit{Computers \& Security}~\cite{deceptionsurvey2024}
(2024) sintetizează cele mai recente tehnici de înșelăciune cibernetică,
incluzând redirecționarea, mutarea țintelor (MTD) și deceptia bazată pe AI.
Lucrarea identifică redirecționarea reactivă ca pe o direcție emergentă cu
puține implementări academice complete o observație care confirmă noutatea
contribuției KernelTrap.
Kahlhofer și Rass~\cite{kahlhofer2024applayer} demonstrează la EuroS\&PW 2024
fezabilitatea unei deceptii la nivel aplicativ fără modificarea aplicațiilor
protejate, principiu reflectat în KernelTrap prin faptul că redirecționarea
atacatorului nu necesită modificări ale daemon-ului SSH sau ale aplicațiilor
existente pe gazda compromisă.

În plus, conceptul de \textbf{honeytoken}în implementarea KernelTrap sub
forma fișierelor canary plasate în containerul honeypot este susținut de
lucrarea seminală a lui Bourke și Grzelak~\cite{bourke2018awstokens} (Black Hat
Asia 2018), care demonstrează empiric o rată de declanșare de 83\% a
token-urilor AWS scurse pe GitHub, și de sondajul recent al lui Khan et
al.~\cite{honeytokenssurvey2024} privind honeytoken-urile în apărarea activă.

\section{Sisteme integrate similare KernelTrap}

Două lucrări academice publicate în 2025 abordează probleme suficient de
similare cu KernelTrap pentru a merita o comparație directă, ele reprezentând
cele două extremități ale spațiului soluțiilor în care se situează contribuția
prezentă.

\textbf{Atawneh et al.}~\cite{atawneh2025symmetry} (2025, \textit{Symmetry})
propun un sistem de detectare a intruziunilor bazat pe Isolation Forest antrenat
pe setul BETH, integrat cu un honeypot pentru colectarea datelor.
Diferențele de design față de KernelTrap sunt următoarele: lucrarea citată
folosește honeypot-ul exclusiv ca sursă de date pentru antrenarea modelului, nu
ca mecanism de răspuns activ; nu implementează agregare temporală cu fereastră
glisantă; nu propune un mecanism de pivotare dinamică; și nu funcționează ca
sistem de producție distribuit cu mai mulți agenți.
Cu toate acestea, valoarea AUROC raportată (0,87) constituie un punct de
referință direct comparabil pentru evaluarea modelului de detecție al
KernelTrap.

\textbf{Satpute et al.}~\cite{satpute2025aidriven} (2025,
\textit{ICCK Transactions on Cybersecurity}) descriu un sistem de detectare a
intruziunilor bazat pe Random Forest și Autoencoder antrenate pe date colectate
dintr-un honeypot SSH (Cowrie).
Abordarea lor este parțial supervizată folosește etichete binare extrase
din log-urile honeypot-ului și nu acoperă scenariul detectării atacatorilor
pe sisteme legitime, ci doar analizează atacatori care au ales să se conecteze
deja la honeypot.
KernelTrap inversează această logică: detectarea are loc pe gazda reală, iar
honeypot-ul este destinația pivotării, nu sursa observațiilor primare.

O lucrare aflată la jumătatea drumului dintre cele două este
\textit{AI-Powered Intrusion Detection System with Honeypot
Integration}~\cite{aipoweredids2025}, care combină detectarea ML cu un honeypot
de tip Cowrie într-un cadru integrat, fără însă a aborda mecanismul de pivotare
automată.
Sistemul rămâne o concatenare de două componente independente, fără bucla de
feedback prezentă în KernelTrap.

\section{Sinteză și poziționarea KernelTrap}

Analiza literaturii recente relevă două observații importante.
Pe de o parte, fiecare componentă individuală a KernelTrap se sprijină pe
lucrări academice solide: monitorizarea prin eBPF/Tracee este validată
de~\cite{soldani2023ebpf,hao2024ebpfruntime,syairozi2025ebpfcomparative,fournier2021ebpf},
detectarea anomaliilor prin Isolation Forest pe BETH este susținută
de~\cite{highnam2021beth,liu2008iforest,atawneh2025symmetry,eremin2025streams},
honeypot-ul cu interacțiune ridicată Cowrie-style este analizat
de~\cite{qcowrie2026,sladic2024llm,honeypotsurvey2025},
iar pivotarea reactivă este pionierată de Lopez
et al.~\cite{lopez2024tcpredirection}.
Pe de altă parte, niciuna dintre lucrările examinate nu integrează aceste patru
componente într-un sistem unitar.
Tabelul~\ref{tab:cap3-comparatie} sintetizează această observație.

\begin{table}[h]
  \centering
  \caption{Comparație între soluțiile existente și KernelTrap pe cele patru
           dimensiuni cheie}
  \label{tab:cap3-comparatie}
  \renewcommand{\arraystretch}{1.08}
  \begin{tabular}{|l|c|c|c|c|}
    \hline
    \textbf{Sistem} & \textbf{eBPF/syscall} & \textbf{ML} &
    \textbf{Honeypot} & \textbf{Răspuns activ} \\
    \hline\hline
    Falco / Tetragon / Tracee                        & Da  & Nu  & Nu      & Nu  \\
    \hline
    Atawneh et al.~\cite{atawneh2025symmetry}        & Da  & Da  & Parțial & Nu  \\
    \hline
    Satpute et al.~\cite{satpute2025aidriven}        & Nu  & Da  & Da      & Nu  \\
    \hline
    AI-Powered IDS~\cite{aipoweredids2025}           & Nu  & Da  & Da      & Nu  \\
    \hline
    Lopez et al.~\cite{lopez2024tcpredirection}      & Nu  & Nu  & Da      & Da  \\
    \hline
    \textbf{KernelTrap}                              & \textbf{Da} & \textbf{Da} &
    \textbf{Da} & \textbf{Da} \\
    \hline
  \end{tabular}
\end{table}

Sistemele comerciale precum Falco, Tetragon, Tracee și KubeArmor oferă
monitorizare excelentă, dar nu includ nici detecție bazată pe ML, nici
componente honeypot, nici răspuns activ prin pivotare.
Lucrările academice apropiate metodologic (Atawneh 2025, Satpute 2025,
AI-Powered IDS 2025) tratează ML-HIDS și honeypot-ul ca componente separate,
fără bucla de feedback care unifică KernelTrap.
Lucrarea lui Lopez et al. demonstrează fezabilitatea pivotării transparente,
dar bazată pe reguli statice de rețea, nu pe decizie ML calibrată
comportamental.

Contribuția specifică a KernelTrap, pe care prezenta lucrare urmărește să o
demonstreze experimental în capitolele următoare, constă în următoarea
combinație: (1) un agent eBPF distribuit, fără stare, cu filtrare strictă pe
sesiuni SSH active și excluderea propriei activități pentru a evita bucla de
auto-întărire, care minimizează zgomotul și suprafața de atac; (2) un model
Isolation Forest extins la 22 de caracteristici (incluzând derivate din
\texttt{args}, \texttt{processName}, \texttt{timestamp} și \texttt{threadId},
ignorate de literatura BETH anterioară), cu praguri adaptive per-(gazdă,
utilizator) pe două niveluri de severitate, antrenat pe BETH cu hiperparametri
ajustați specific scenariului syscall (\texttt{n\_estimators}~= 2000, percentile
2,0\% și 0,2\%) și complementat de o listă albă de procese de infrastructură
pentru atenuarea derivei de distribuție; (3) un mecanism de agregare prin
fereastră glisantă cu prag dublu (cantitativ și de rată) care convertește scoruri
zgomotoase în decizii robuste; (4) o pivotare transparentă declanșată automat
care combină semnale POSIX, blocare \texttt{iptables} a conexiunilor noi și
\texttt{exec} către un container honeypot Docker cu identitate oglindită,
fără reguli NAT și fără reconectarea sesiunii SSH; și (5) un panou de control
în timp real care permite și intervenție manuală.
Această combinație nu apare în nicio lucrare publicată cunoscută și constituie
principala contribuție a tezei de față.
